\section*{Chapter \ref{ch:nw:colocation-products-and-how-they-are-sold} --- Colocation Products and How They Are Sold}

\begin{solution}{exo:nw:colocation-products-and-how-they-are-sold:1}
$(120 - 30) \times 4.88 \approx \qty{439}{\nano\second}$.
\end{solution}

\begin{solution}{exo:nw:colocation-products-and-how-they-are-sold:2}
$7\,230/10 = 723$~USD per kilowatt a month at \qty{10}{\kilo\watt}, $7\,230/15 = 482$~USD at \qty{15}{\kilo\watt}: the filing's own figures.
\end{solution}

\begin{solution}{exo:nw:colocation-products-and-how-they-are-sold:3}
It raised it by $2 \times 12 \times 1\,500 = 36\,000$~USD, from 324\,000 to 360\,000~USD a year.
\end{solution}

\begin{solution}{exo:nw:colocation-products-and-how-they-are-sold:4}
Power, and the cooling that removes it, are the building's scarce and expensive resources; a cabinet's cost to the operator is mostly the power
and cooling capacity it reserves, which the \omterm{def:nw:servers:ru}{rack units} do not measure.
\end{solution}

\begin{solution}{exo:nw:colocation-products-and-how-they-are-sold:5}
They forbid a direct \omterm{def:nw:colocation-products-and-how-they-are-sold:mmr}{cross-connect} between participants that could result in a trade inside the building, and require trading on an authorised
venue to go through the venue's network. Venues write such rules so that trading happens on regulated venues, under their surveillance and
fairness rules, and so that no private link inside the building bypasses them.
\end{solution}

\begin{solution}{exo:nw:colocation-products-and-how-they-are-sold:6}
Faster servers and code inside the cabinet; faster devices and fewer hops inside the firm's own cage; faster links between buildings for
firms trading across venues (and faster venue connection products, where a venue sells more than one class).
\end{solution}

\begin{solution}{exo:nw:colocation-products-and-how-they-are-sold:7}
368\,053~USD a year over 36 months, 359\,648~USD over 60: 8\,405~USD less, because the 63\,040~USD of one-time fees weigh less per month.
\end{solution}

\begin{solution}{exo:nw:colocation-products-and-how-they-are-sold:8}
Equal cable lengths equalise the venue's side only: each firm's own switches, cards, servers and software still differ, and firms taking
different products (a standard or a lower-latency network, different port speeds) are equal only within each product.
\end{solution}

\begin{solution}{pb:nw:colocation-products-and-how-they-are-sold:1}
\begin{enumerate}
\item $2 \times 7\,230 = 14\,460$~USD.
\item $2 \times (5\,940 + 4\,560) = 21\,000$~USD, about 583~USD a month over 36 months.
\item $2 \times 15\,000 = 30\,000$~USD.
\item $14\,460 + 583 + 30\,000 \approx 45\,043$~USD a month, 540\,520~USD a year.
\item $50\,000 \times 20 = 1\,000\,000$ contracts.
\item About 0.045~USD per contract.
\item About 0.09~USD per contract: the footprint's cost is fixed.
\item The connections: 30\,000 of the 45\,043~USD; one connection instead of two, if the firm accepts the loss of redundancy.
\item $(150 - 10) \times 4.88 \approx \qty{683}{\nano\second}$.
\item The venue cuts every connection to the same length; buying a shorter one would give it an advantage the fairness rules forbid.
\item That the same conditions (space, power, cooling, cable length, connectivity) apply to all users of the same service, that latency is
  monitored, and that services are sold unbundled.
\item The venue's lower-latency network and a 10-gigabit port for bursts (chapter~1), as a product sold equally to all who buy it.
\item \textbf{Named result.} About 45\,043~USD a month all in, 0.045~USD per contract at 50\,000 contracts a day.
\item $21\,000 / (45\,043 \times 36)$: 1.3\%.
\item At least every 60 days and after every fee filing: fees change by filing, and U.S. exchanges file often.
\item The firm's own equipment, market-data fees, \omterm{def:nw:colocation-products-and-how-they-are-sold:mmr}{cross-connects} to carriers, \omterm{def:nw:colocation-products-and-how-they-are-sold:hands}{remote hands}, and the disaster-recovery site's footprint
  (chapter~28).
\item The total would shift toward space and power; the break-even rises or falls with the sum, not with any line.
\item Connectivity into the venue is a product only the venue sells, priced by filing on the value of speed; space and power are commodities.
\item Volume-linked tiers and incentive programmes with the venue, carrier and \omterm{def:nw:colocation-products-and-how-they-are-sold:mmr}{cross-connect} prices with the data-centre operator, and the
  term of the contract.
\item Proximity on equal terms: the right to be in the building, on the same cable as everyone else, at a price.
\end{enumerate}
\end{solution}

\begin{solution}{iq:nw:colocation-products-and-how-they-are-sold:1}
Space for equipment in the venue's data centre, power and cooling for it, \omterm{def:nw:colocation-products-and-how-they-are-sold:mmr}{cross-connects} to the venue's network and to carriers, the venue's
market-data and order-entry ports, and services such as \omterm{def:nw:colocation-products-and-how-they-are-sold:hands}{remote hands}; most of the cost is power and the venue's connectivity products.
\iqlookfor{power and connectivity as the real costs.}
\end{solution}

\begin{solution}{iq:nw:colocation-products-and-how-they-are-sold:2}
A \omterm{def:nw:colocation-products-and-how-they-are-sold:mmr}{cross-connect} is an operator-installed cable between two points of the building; the \omterm{def:nw:colocation-products-and-how-they-are-sold:mmr}{meet-me room} is where carriers terminate. A carrier's
circuit ends on its equipment in the \omterm{def:nw:colocation-products-and-how-they-are-sold:mmr}{meet-me room}, and a \omterm{def:nw:colocation-products-and-how-they-are-sold:mmr}{cross-connect} (often through a patch panel and a distribution point) carries it to
the customer's cabinet.
\iqlookfor{the physical path and who manages it.}
\end{solution}

\begin{solution}{iq:nw:colocation-products-and-how-they-are-sold:3}
To give every customer of the same service the same latency to the venue, as fairness rules require. What remains is everything outside the
venue's cable: the firm's own hardware and software, its choice of connectivity products, and its links between venues.
\iqlookfor{equal on the venue's side, competition everywhere else.}
\end{solution}

\begin{solution}{iq:nw:colocation-products-and-how-they-are-sold:4}
Power and cabinet space first (long lead times), then \omterm{def:nw:colocation-products-and-how-they-are-sold:mmr}{cross-connects} and the venue's ports, the carriers' circuits for the firm's own links,
timing (an antenna feed or the operator's service), then the equipment, installation by \omterm{def:nw:colocation-products-and-how-they-are-sold:hands}{remote hands} or the firm's staff, and the venue's
certification of the new connections.
\iqlookfor{lead times and dependencies in the right order.}
\end{solution}

\begin{solution}{iq:nw:colocation-products-and-how-they-are-sold:5}
Price the footprint all in, from the current schedules, with one-time fees amortised; compare it with the strategy's expected net edge per unit
times its volume at that venue, with and without the latency the footprint buys; include the cost of its disaster-recovery twin.
\iqlookfor{a break-even per unit traded, and the counterfactual without \omterm{def:nw:colocation-products-and-how-they-are-sold:colo}{colocation}.}
\end{solution}

\begin{solution}{iq:nw:colocation-products-and-how-they-are-sold:6}
In the exchange's fee schedule in its rulebook, and in the rule filings that change it, published by the SEC and in the Federal Register;
check the filing dates and re-read after every filing.
\iqlookfor{public filings and a date on every figure.}
\end{solution}
